\documentclass{situationist-preprint}

\title{Self-Formation Without Regress}
\msid{ACT-08-b}
\msfamily{ACT-08}
\msdiscipline{Action, Agency and Free Will}
\mstype{Companion manuscript}
\mskeywords{Basic Argument; self-formation; ultimate responsibility; character; Strawson; tracing}

\begin{document}
\maketitle

\begin{abstract}
ACT-08 argues that Strawson's Basic Argument fails at its ultimacy premise, and that responsibility for how one is can be acquired by degrees, as responsible actions shape the nature from which later actions flow. We develop the second claim. We show that responsibility for one's nature, so understood, accumulates without regress, because each step requires responsibility for an action, not for the whole nature from which it issued. We then answer the objection that an unchosen starting point taints everything built on it. Self-formation begins from a nature one did not choose, and it need not begin from nothing.
\end{abstract}

\section{Graded self-formation}

ACT-08-a accepts Strawson's regress premise as an analysis of ultimate responsibility (Principle~P3) and grants that no one is ultimately responsible for her nature (Lemma~3.1). It rejects the premise that responsibility for action requires such responsibility (Principle~P2), on the ground that action from an owned, reasons-responsive mechanism suffices (Assumption~2.4); and it shows that, given that actions contribute to later natures (Assumption~2.6), every later nature is partly formed by actions for which the agent is responsible (Proposition~3.3). We now say what responsibility for one's nature comes to.

\begin{definition}\label{d:share}
The \emph{self-formed share} of an agent's nature at stage $n$ is the part of it that is due to actions for which she was responsible at earlier stages.
\end{definition}

\begin{proposition}\label{p:accum}
Given the assumptions of ACT-08-a, with $k$ the first stage of Assumption~2.5 there, if each responsible action contributes to the next nature and the contributions persist, the self-formed share of the agent's nature is zero at stages up to $k$ and positive at every stage after $k$, and it does not decrease so long as she continues to act responsibly.
\end{proposition}

\begin{proof}
Before stage $k$ the agent is not responsible for her actions since $k$ is the first stage at which she acts from an owned, reasons-responsive mechanism, so no part of her nature up to $k$ is due to responsible action. By Proposition~3.3 of ACT-08-a, every later nature is partly formed by a responsible action, so its share is positive. If contributions persist, each later nature retains what earlier responsible actions contributed and adds what the latest one contributes.
\end{proof}

No regress arises, because each increment requires only that the agent was responsible for an action, and responsibility for an action requires, by Assumption~2.4 of ACT-08-a, only that it issue from an owned, reasons-responsive mechanism. It does not require prior responsibility for the whole nature from which the action issued. This is the structure of Aristotle's account of character, on which we are in part responsible for our states because we are responsible for the activities that produced them \citep[III.5]{AristotleNE}, and of the tracing accounts that ground responsibility for a present state in responsibility for earlier conduct \citep{FischerRavizza1998}.

\section{The tainted origin}

\begin{objection}[Inherited non-responsibility]
The first responsible action issued from an unchosen nature. Everything that it contributed was produced by a nature for which the agent was not responsible, so the contribution is not hers in any deep sense, and the same holds at every later step \citep{Strawson1994}.
\end{objection}

\begin{reply}
The objection assumes that responsibility for a product requires responsibility for every factor that produced it. That is Principle~P2 of ACT-08-a applied to natures instead of actions, and it is exactly the premise in dispute. An agent who resolves to be more patient, and becomes so through practice, is responsible for her patience because she is responsible for the resolution and the practice; that her capacity to resolve was given rather than chosen no more cancels this than the unchosen strength of her arms cancels her responsibility for lifting a child out of danger. Responsibility is not inherited from origins in either direction: unchosen origins neither confer nor withhold it. This is the point of \citet{Dennett1984}'s account of self-made selves: the self-creation worth wanting is the gradual kind that finite agents can achieve.
\end{reply}

\section{Conclusion}

Self-formation accumulates by steps, each requiring responsibility for an action and none requiring responsibility for one's whole prior self. It therefore starts from an unchosen nature without starting from nothing, and without regress. This supports the conclusion of ACT-08: we need not be the ultimate source of our actions, and the Basic Argument fails because self-formation need not start from nothing.

\begin{thebibliography}{9}
\bibitem[Aristotle(NE)]{AristotleNE} Aristotle. \emph{Nicomachean Ethics}. Cited by book and chapter.
\bibitem[Dennett(1984)]{Dennett1984} Dennett, D.~C. (1984). \emph{Elbow Room: The Varieties of Free Will Worth Wanting}. MIT Press.
\bibitem[Fischer and Ravizza(1998)]{FischerRavizza1998} Fischer, J.~M. and Ravizza, M. (1998). \emph{Responsibility and Control: A Theory of Moral Responsibility}. Cambridge University Press.
\bibitem[Strawson(1994)]{Strawson1994} Strawson, G. (1994). The impossibility of moral responsibility. \emph{Philosophical Studies}, 75, 5--24.
\end{thebibliography}

\end{document}
